\documentclass{eic-bcc-tcc}

\title{Metaverso: uma análise comparativa de definições, tecnologias e aplicações}
\titleen{Metaverse: a Comparative Analysis of Definitions, Technologies and Applications}

% Preencha estes campos com os dados reais do trabalho.
\addstudent[Full Name]{Nome Completo}
\advisor[f]{Nome da Orientadora, D.Sc.}
\date{2026}

\newacronym{ia}{IA}{Inteligência Artificial}
\newacronym{xr}{XR}{Realidade Estendida}
\newacronym{vr}{VR}{Realidade Virtual}
\newacronym{ar}{AR}{Realidade Aumentada}
\newacronym{mr}{MR}{Realidade Mista}
\newacronym{prisma}{PRISMA}{Preferred Reporting Items for Systematic Reviews and Meta-Analyses}

\begin{document}

\maketitle
\makeapproval

\begin{resumo}
Este trabalho apresenta uma análise comparativa do conceito de metaverso a partir da definição adotada no projeto e de trabalhos científicos de revisão da literatura. O estudo confronta uma definição orientada à experiência do usuário com a perspectiva tecnológica apresentada por Tukur et al. (2024), além de utilizar a revisão sistemática de Ritterbusch e Teichmann (2023) como apoio conceitual. A análise separa características do fenômeno, propriedades desejáveis e tecnologias habilitadoras, evitando o tratamento de tecnologias específicas como requisitos universais. Como resultado, propõe-se uma definição operacional em camadas, incorporando ambientes tridimensionais compartilhados, interação por identidades digitais, imersão, persistência, interoperabilidade e economia digital como dimensões do ecossistema, sustentadas por tecnologias variáveis. O trabalho também estrutura um procedimento de busca baseado nas recomendações do PRISMA 2020 e apresenta um quadro de registro da busca exploratória, deixando explícita a diferença entre essa etapa e uma revisão sistemática completa.
\palavraschave{Metaverso. Revisão da literatura. PRISMA. Realidade estendida. Tecnologias digitais.}
\end{resumo}

\begin{abstract}
This study presents a comparative analysis of the metaverse concept based on the definition adopted in the project and scientific literature reviews. The analysis contrasts a user-experience-oriented definition with the technological perspective presented by Tukur et al. (2024), while also using the systematic review by Ritterbusch and Teichmann (2023) as conceptual support. The study separates properties of the phenomenon, desirable characteristics, and enabling technologies, avoiding the treatment of specific technologies as universal requirements. As a result, an operational layered definition is proposed, incorporating shared three-dimensional environments, interaction through digital identities, immersion, persistence, interoperability, and digital economy as dimensions of the ecosystem, supported by variable technologies. The work also structures a search procedure based on PRISMA 2020 recommendations and presents a record of the exploratory search, explicitly distinguishing this stage from a complete systematic review.
\keywords{Metaverse. Literature review. PRISMA. Extended reality. Digital technologies.}
\end{abstract}

\listoftablespage
\listofacronymspage
\tableofcontentspage

\mainmatter

\chapter{Introdução}
\label{cha:introducao}

O metaverso é um conceito multidimensional cuja definição permanece objeto de debate na literatura científica. Revisões recentes tratam o tema por perspectivas diferentes: algumas concentram-se na definição conceitual; outras investigam tecnologias, aplicações, arquiteturas ou interoperabilidade. Essa diversidade torna inadequado tratar uma única formulação como definição universal.

Ritterbusch e Teichmann (2023) realizaram uma revisão sistemática especificamente dedicada à definição do metaverso e analisaram 28 definições e descrições científicas. O estudo mostra a variedade de formulações utilizadas na literatura e constitui uma referência para a análise conceitual realizada neste trabalho. \cite{ritterbusch2023}

A literatura posterior amplia o problema. Tukur et al. (2024) realizaram uma revisão de escopo sobre técnicas, tecnologias e aplicações do metaverso, identificando 12 classes tecnológicas e aplicações em 11 setores. Boztas, Ghadafi e Ibrahim (2025), por sua vez, realizaram uma revisão sistemática de 103 estudos, dos quais 33 apresentavam definições explícitas e 25 propunham modelos arquiteturais. \cite{tukur2024,boztas2025}

Diante desse cenário, o presente trabalho não pretende declarar uma nova definição universal. Seu objetivo é comparar a definição adotada no projeto com perspectivas científicas verificáveis e, a partir dessa comparação, construir uma definição operacional adequada ao escopo do projeto.

A questão central é, portanto, distinguir três níveis que frequentemente aparecem misturados: as propriedades que caracterizam o fenômeno, as tecnologias que podem sustentá-lo e as aplicações que podem ser desenvolvidas sobre ele. Essa separação é necessária para evitar que tecnologias específicas, como realidade virtual, inteligência artificial ou blockchain, sejam tratadas automaticamente como requisitos universais.

O objetivo geral é analisar comparativamente a literatura científica sobre o metaverso e confrontá-la com a definição de trabalho adotada no projeto. Como objetivos específicos, busca-se: identificar características recorrentes nas definições; comparar definições e arquiteturas recentes; distinguir propriedades conceituais de tecnologias habilitadoras; analisar aplicações e interoperabilidade; e formular uma definição operacional para o projeto.

A contribuição deste trabalho é, portanto, uma síntese comparativa fundamentada em fontes científicas identificáveis. Sempre que um número é apresentado, ele é atribuído à fonte que o produziu. Não são apresentados como resultados desta pesquisa números que não tenham sido efetivamente coletados e registrados pelo presente estudo.

O restante do trabalho está organizado da seguinte forma: o Capítulo~\ref{cha:revisao} apresenta a revisão da literatura; o Capítulo~\ref{cha:questoes} apresenta as questões de pesquisa; o Capítulo~\ref{cha:relacionados} discute os trabalhos relacionados; o Capítulo~\ref{cha:metodo} descreve o método e a rastreabilidade das evidências; o Capítulo~\ref{cha:critica} apresenta a análise crítica e a definição operacional proposta; o Capítulo~\ref{cha:resultados} sintetiza os resultados; e o Capítulo~\ref{cha:conclusao} apresenta as conclusões e limitações.

\chapter{Revisão da Literatura}
\label{cha:revisao}

\section{Metaverso}
\label{sec:metaverso}

O termo metaverso não possui uma única definição consolidada. Ritterbusch e Teichmann (2023), em uma revisão sistemática dedicada especificamente à definição do conceito, identificaram diferentes formulações na literatura e destacaram atributos como ambientes tridimensionais, representação dos usuários por avatares, interação, imersão, economia virtual e tecnologias como realidade virtual, realidade aumentada e blockchain.

A definição adotada neste trabalho acrescenta à dimensão tridimensional e imersiva três características particularmente relevantes para o escopo do projeto: persistência, interoperabilidade e economia digital. A persistência corresponde à continuidade do ambiente mesmo quando determinado usuário deixa a plataforma. A interoperabilidade representa a possibilidade, ainda conceitual ou parcial, de transportar identidades, objetos ou outros elementos entre diferentes ambientes. A economia digital corresponde à existência de mecanismos de troca, propriedade, produção ou consumo de bens e serviços digitais.

A literatura mais recente reforça que o metaverso deve ser compreendido como um ecossistema tecnológico, e não somente como um ambiente visual tridimensional. Tukur et al. (2024) realizaram uma revisão de escopo de trabalhos publicados entre 2020 e 2022 e identificaram quatro técnicas principais e doze classes tecnológicas relacionadas à criação de ambientes digitais para o metaverso. Entre as tecnologias mais recorrentes estavam a realidade estendida, a inteligência artificial e tecnologias descentralizadas.

O estudo de Tukur et al. (2024) também evidencia que as aplicações do metaverso ultrapassam o entretenimento. Os autores identificaram aplicações em onze setores, com destaque para educação, manufatura, saúde e mercado imobiliário. Esse resultado é relevante porque desloca a discussão de uma visão exclusivamente associada a jogos para uma perspectiva de infraestrutura digital e interação entre pessoas, objetos e serviços.

Uddin et al. (2024) ampliam essa perspectiva ao analisar a convergência entre metaverso, inteligência artificial e blockchain. O estudo discute como essas tecnologias podem atuar conjuntamente na criação de ambientes virtuais, na personalização de experiências, na gestão de ativos digitais e em questões de segurança. Essa abordagem ajuda a separar o conceito de metaverso das tecnologias que podem viabilizá-lo: realidade virtual, inteligência artificial e blockchain são componentes possíveis, mas não devem ser tratados automaticamente como sinônimos do conceito.

Assim, a revisão indica uma distinção importante. Algumas propriedades aparecem como características conceituais recorrentes, como ambientes digitais tridimensionais, interação por avatares e experiências compartilhadas. Outras propriedades, como blockchain, inteligência artificial, realidade virtual e economias baseadas em ativos digitais, funcionam principalmente como tecnologias ou mecanismos de implementação. A definição adotada neste trabalho deve preservar essa distinção.

\section{Questões da Pesquisa}
\label{cha:questoes}

Para orientar a análise comparativa, foram formuladas as seguintes questões de pesquisa:

\begin{itemize}
    \item \textbf{QP1:} Quais características aparecem com maior recorrência nas definições científicas de metaverso?
    \item \textbf{QP2:} Quais características da definição adotada neste projeto são sustentadas pela literatura analisada?
    \item \textbf{QP3:} Quais características podem ser unificadas sem ampliar excessivamente o conceito de metaverso?
    \item \textbf{QP4:} Quais características devem ser tratadas como possibilidades de implementação, e não como requisitos necessários?
    \item \textbf{QP5:} Quais lacunas permanecem na definição adotada e quais consequências essas lacunas podem trazer para o desenvolvimento do projeto?
\end{itemize}

As questões foram formuladas para evitar que a comparação se limite a uma escolha entre duas definições. O objetivo é identificar convergências, divergências e níveis diferentes de abstração. Dessa maneira, uma característica pode ser mantida na definição final como elemento conceitual, enquanto outra pode ser deslocada para a camada tecnológica ou para a descrição das aplicações.

\section{Trabalhos Relacionados}
\label{cha:relacionados}

O primeiro trabalho utilizado para comparação é o estudo de Ritterbusch e Teichmann (2023), que realiza uma revisão sistemática especificamente voltada à definição do metaverso. Sua principal contribuição para este projeto é conceitual: o estudo demonstra que não existe consenso absoluto e organiza atributos recorrentes das definições encontradas.

O segundo trabalho, utilizado como referência principal para a comparação aprofundada, é a revisão de escopo de Tukur et al. (2024). Diferentemente do foco predominantemente conceitual de Ritterbusch e Teichmann, Tukur et al. concentram-se nas técnicas, tecnologias e aplicações necessárias à criação de ambientes digitais associados ao metaverso. O estudo utilizou procedimentos de revisão baseados no PRISMA-ScR e pesquisou ACM, IEEE, Scopus e Google Scholar, mantendo trinta trabalhos primários considerados relevantes.

Para este projeto, a comparação central é entre a definição de trabalho adotada e a perspectiva apresentada por Tukur et al. (2024). A definição do projeto enfatiza a experiência do usuário e as propriedades esperadas do ambiente; Tukur et al. enfatizam a infraestrutura tecnológica e os setores de aplicação. Portanto, os dois enfoques não são concorrentes: eles descrevem dimensões diferentes do mesmo objeto.

\begin{table}[H]
\centering
\caption{Comparação entre a definição do projeto e o estudo de Tukur et al. (2024)}
\label{tab:comparacao}
\begin{tabular}{p{0.24\textwidth}p{0.33\textwidth}p{0.33\textwidth}}
\toprule
\textbf{Dimensão} & \textbf{Definição do projeto} & \textbf{Tukur et al. (2024)}\\
\midrule
Ambiente tridimensional & Elemento central da definição & Relacionado à criação de ambientes virtuais\\
Imersão & Destacada na experiência do usuário & Associada principalmente à XR e à convergência de VR\\
Avatares & Meio de representação e interação & Compatível com a interação em ambientes virtuais\\
Persistência & Considerada característica desejável do ambiente & Não aparece como eixo principal da taxonomia tecnológica\\
Interoperabilidade & Considerada como possibilidade estrutural & Relacionada aos desafios de integração e conexão entre tecnologias\\
Economia digital & Considerada parte da experiência ampliada & Aparece como dimensão de aplicação e como consequência do ecossistema tecnológico\\
IA & Tecnologia de suporte & Uma das tecnologias mais recorrentes\\
Blockchain/descentralização & Tecnologia de suporte à economia e aos ativos & Uma das tecnologias recorrentes\\
Aplicações & Reuniões, aulas, shows, compras e experiências digitais & Onze setores identificados na revisão\\
\bottomrule
\end{tabular}
\end{table}

A comparação mostra que não há necessidade de escolher entre as duas perspectivas. A definição do projeto pode funcionar como uma camada conceitual orientada à experiência e às propriedades do ambiente, enquanto a revisão de Tukur et al. fornece a camada tecnológica necessária para explicar como tais propriedades podem ser implementadas. Essa complementaridade fundamenta a proposta de unificação apresentada posteriormente.

\chapter{Método}
\label{cha:metodo}

\section{Delineamento da pesquisa}

Este estudo adota uma revisão bibliográfica estruturada, de caráter qualitativo e comparativo, orientada pelas recomendações do PRISMA 2020 e pela extensão PRISMA-S para relato das estratégias de busca. O objetivo é identificar como a literatura científica recente define o metaverso, quais características são recorrentes, quais tecnologias aparecem como habilitadoras e quais aplicações são descritas. A revisão é utilizada como base para confrontar a definição adotada no projeto e construir uma definição operacional em camadas.

O PRISMA 2020 orienta o relato transparente das fontes de informação, critérios de elegibilidade, processo de seleção e síntese. O PRISMA-S acrescenta requisitos específicos para que a busca seja reproduzível, incluindo o registro da fonte, da data e da estratégia completa utilizada em cada fonte \cite{page2021,rethlefsen2021}.

\section{Questões de pesquisa}

A busca foi estruturada para responder às seguintes questões:

\begin{itemize}
    \item \textbf{QP1:} Quais características aparecem com maior recorrência nas definições científicas de metaverso?
    \item \textbf{QP2:} Quais características da definição adotada neste projeto são sustentadas pela literatura?
    \item \textbf{QP3:} Quais características podem ser unificadas sem ampliar excessivamente o conceito?
    \item \textbf{QP4:} Quais elementos devem ser tratados como tecnologias habilitadoras, e não como requisitos conceituais?
    \item \textbf{QP5:} Quais lacunas permanecem na definição adotada e quais implicações elas trazem para o projeto?
\end{itemize}

\section{Período e fontes de informação}

O protocolo foi definido para contemplar publicações em inglês entre janeiro de 2022 e setembro de 2026, período escolhido para concentrar a análise na literatura produzida após a consolidação do debate acadêmico contemporâneo sobre o metaverso.

As fontes planejadas são:

\begin{itemize}
    \item Scopus;
    \item Web of Science;
    \item IEEE Xplore;
    \item ACM Digital Library;
    \item Google Scholar, como fonte complementar de recuperação.
\end{itemize}

A escolha é coerente com revisões anteriores do tema. Tukur et al. utilizaram IEEE Xplore, Scopus, ACM e Google Scholar e reportaram 447 registros após restrições de busca; Ritterbusch e Teichmann utilizaram Web of Science, JSTOR, Wiley e AIS, com 381 registros iniciais; e Boztas et al. realizaram uma revisão de definições e arquiteturas a partir do Google Scholar, iniciando com 237 registros \cite{tukur2024,ritterbusch2023,boztas2025}. Esses números pertencem às respectivas revisões e não são tratados como resultados desta pesquisa.

\section{Estratégia de busca}

Foi definida uma string conceitual principal:

\begin{verbatim}
("metaverse" OR "metaverses")
AND
(definition OR conceptualization OR characteristics OR architecture)
AND
(technology OR technologies OR application OR applications
OR interoperability OR persistence)
\end{verbatim}

Também foi definida uma string específica para localização de revisões:

\begin{verbatim}
("metaverse" OR "metaverses")
AND
("systematic literature review" OR "systematic review"
OR "scoping review" OR "literature review")
AND
(definition OR characteristics OR technologies OR applications)
\end{verbatim}

A sintaxe deve ser adaptada aos operadores e campos específicos de cada base, mantendo a mesma lógica conceitual. De acordo com o PRISMA-S, a versão efetivamente executada em cada base deve ser preservada exatamente como utilizada, juntamente com a data e os filtros aplicados \cite{rethlefsen2021}.

\section{Critérios de inclusão}

Foram definidos previamente os seguintes critérios:

\begin{itemize}
    \item publicação entre janeiro de 2022 e setembro de 2026;
    \item publicação em inglês;
    \item artigo científico ou trabalho de conferência com contribuição acadêmica identificável;
    \item relação direta com o conceito de metaverso;
    \item discussão de pelo menos uma das dimensões de interesse: definição, características, arquitetura, tecnologias, aplicações, persistência ou interoperabilidade;
    \item texto completo disponível para avaliação;
    \item contribuição suficiente para responder pelo menos uma questão de pesquisa.
\end{itemize}

\section{Critérios de exclusão}

Foram definidos os seguintes critérios de exclusão:

\begin{itemize}
    \item notícias, páginas comerciais, textos de opinião e materiais sem caráter científico;
    \item trabalhos que apenas mencionem o termo metaverso sem analisá-lo;
    \item registros duplicados;
    \item trabalhos sem texto completo disponível;
    \item trabalhos sem relação suficiente com as questões de pesquisa;
    \item estudos cujo foco seja exclusivamente uma tecnologia isolada, sem relação analítica com o metaverso;
    \item documentos fora do período ou idioma definidos.
\end{itemize}

\section{Execução e auditoria da busca}

A busca bibliográfica foi executada em 29 de setembro de 2026 por meio das fontes científicas e páginas públicas acessíveis durante a pesquisa. Foram localizados e verificados registros diretamente relacionados à definição, às características, às arquiteturas, às tecnologias e às aplicações do metaverso.

Entretanto, os contadores ao vivo de resultados das interfaces fechadas de Scopus, Web of Science, IEEE Xplore, ACM Digital Library e Google Scholar não são expostos de maneira confiável no ambiente utilizado nesta pesquisa. Por essa razão, este trabalho não atribui números inventados a essas bases. Os quantitativos finais do fluxograma PRISMA somente poderão ser declarados como resultados desta revisão após a execução das strings nas interfaces das bases e o registro dos respectivos contadores.

Essa decisão metodológica é deliberada: um número de resultados não observado diretamente não deve ser apresentado como dado empírico da revisão.

Como controle externo de plausibilidade, foram comparados os procedimentos e quantitativos publicados por revisões anteriores. Tukur et al. reportaram 212 registros provenientes do IEEE Xplore, 96 do Scopus, 39 da ACM e 100 do Google Scholar, totalizando 447 registros após restrições específicas; 141 duplicatas foram removidas e 122 estudos chegaram à avaliação de texto completo \cite{tukur2024}. Ritterbusch e Teichmann reportaram 381 registros iniciais em Web of Science, JSTOR, Wiley e AIS, chegando a 28 definições ou descrições após as etapas de seleção e buscas de citações \cite{ritterbusch2023}. Boztas et al. identificaram inicialmente 237 registros no Google Scholar, removeram 6 duplicatas e 7 registros por outros motivos, avaliaram 192 textos completos e incluíram 103 estudos \cite{boztas2025}.

\section{Resultados bibliográficos verificados}

A busca realizada permitiu identificar um núcleo atualizado de trabalhos particularmente relevantes para este estudo:

\begin{table}[H]
\centering
\caption{Trabalhos de revisão e definição identificados na busca bibliográfica}
\label{tab:trabalhos_verificados}
\begin{tabular}{p{0.28\textwidth}p{0.23\textwidth}p{0.39\textwidth}}
\toprule
\textbf{Trabalho} & \textbf{Ano} & \textbf{Contribuição para este estudo}\\
\midrule
Ritterbusch e Teichmann & 2023 & Definições e características recorrentes do metaverso.\\
Tukur et al. & 2024 & Técnicas, tecnologias e aplicações do metaverso.\\
Zhou, Chen e Jin & 2024 & Definição, características técnicas e comportamento do usuário em perspectiva sociotécnica.\\
Bénaben, Congès e Fertier & 2025 & Proposta de definição unificada e framework para propriedades do metaverso.\\
Enamorado-Díaz et al. & 2025 & Aplicações, tecnologias e desafios de engenharia de software.\\
Boztas, Ghadafi e Ibrahim & 2025 & Definições e arquiteturas, com análise temática de 33 definições.\\
\bottomrule
\end{tabular}
\end{table}

Os trabalhos identificados reforçam uma convergência importante para o objetivo deste artigo: não há uma definição única e universalmente aceita, mas existem conjuntos recorrentes de características, tecnologias habilitadoras e dimensões sociais, econômicas e arquitetônicas \cite{ritterbusch2023,zhou2024,boztas2025}.

\section{Fluxo PRISMA}

O fluxograma PRISMA definitivo não é preenchido com estimativas. Nesta versão, a tabela abaixo funciona como registro metodológico do estado da busca e separa claramente dados efetivamente observados de dados que ainda dependem da execução manual nas interfaces das bases.

\begin{table}[H]
\centering
\caption{Estado do registro PRISMA}
\label{tab:prisma}
\begin{tabular}{p{0.62\textwidth}p{0.23\textwidth}}
\toprule
\textbf{Etapa} & \textbf{Situação}\\
\midrule
Bases e fontes definidas & Registrado\\
Strings conceituais definidas & Registrado\\
Critérios de inclusão definidos & Registrado\\
Critérios de exclusão definidos & Registrado\\
Busca bibliográfica exploratória executada & Registrado\\
Registros recuperados em cada base fechada & A registrar na interface da base\\
Duplicatas & A calcular após exportação\\
Triagem de títulos e resumos & A executar sobre o conjunto exportado\\
Elegibilidade por texto completo & A executar sobre o conjunto exportado\\
Estudos incluídos & A calcular após a elegibilidade\\
\bottomrule
\end{tabular}
\end{table}

Portanto, esta versão não apresenta o quadro como um PRISMA numérico completo. O preenchimento definitivo depende dos contadores e dos registros exportados das bases. Essa distinção preserva a rastreabilidade e evita que números de revisões anteriores sejam confundidos com resultados desta pesquisa.

\section{Procedimento de análise}

Após a seleção definitiva, cada estudo será classificado segundo cinco dimensões: (i) definição do metaverso; (ii) características estruturais; (iii) tecnologias habilitadoras; (iv) aplicações e setores; e (v) desafios ou limitações.

A síntese será realizada por comparação temática. Características conceituais, como ambiente compartilhado, interação, identidade digital, persistência e imersão, serão separadas das tecnologias que podem viabilizá-las, como realidade estendida, inteligência artificial, computação distribuída e tecnologias descentralizadas. Essa separação evita transformar uma tecnologia específica em requisito universal do conceito.

A etapa final consistirá na comparação entre a definição adotada no projeto e os padrões encontrados na literatura. A definição resultante será apresentada como uma definição operacional para o projeto, e não como uma afirmação de que existe uma única definição científica definitiva.

\chapter{Análise Crítica e Unificação}
\label{cha:critica}

\section{Crítica do trabalho atual}

A análise do projeto disponível no repositório revela que a principal limitação da versão anterior não estava na escolha do tema, mas na ausência de uma delimitação científica suficientemente operacional. O arquivo principal ainda apresentava grande quantidade de conteúdo genérico do template de TCC, enquanto a discussão específica sobre metaverso, critérios de seleção de literatura, comparação de estudos e perguntas de pesquisa ainda não estava consolidada.

Outro ponto crítico é a tendência de reunir em uma única definição elementos de naturezas diferentes. Persistência, interação por avatares e existência de um ambiente tridimensional descrevem propriedades do fenômeno. Já realidade virtual, inteligência artificial, blockchain e computação em nuvem descrevem tecnologias que podem viabilizar essas propriedades. Misturar essas camadas pode tornar a definição circular ou excessivamente dependente de tecnologias específicas.

A terceira limitação é metodológica. Uma tabela PRISMA só é cientificamente válida como registro quantitativo quando os números correspondem a uma busca realmente executada e documentada. Por isso, nesta versão os números apresentados são explicitamente identificados como pertencentes à busca exploratória comparativa, e não como uma revisão sistemática completa.

\section{Comparação conceitual}

A definição do projeto e o estudo de Tukur et al. (2024) apresentam pontos de convergência, mas partem de objetivos diferentes. O projeto procura caracterizar o que o usuário vivencia e quais propriedades o ambiente deve apresentar. Tukur et al. procuram identificar as técnicas e tecnologias que sustentam a construção desses ambientes.

Ritterbusch e Teichmann (2023) fornecem uma terceira camada de validação conceitual ao mostrar que características como ambientes tridimensionais, avatares, interação, imersão e economia virtual aparecem de forma recorrente na literatura, embora não exista consenso absoluto sobre uma única definição.

Dessa comparação, resulta uma conclusão metodológica: a unificação é adequada desde que seja feita por camadas. Não se recomenda simplesmente juntar todas as características em uma lista única e tratá-las como requisitos obrigatórios.

\section{Definição unificada}

Com base na comparação realizada, este trabalho adota a seguinte definição operacional:

\begin{quote}
\textit{O metaverso é um ecossistema de ambientes digitais tridimensionais e compartilhados, nos quais usuários representados por identidades digitais podem interagir entre si e com objetos virtuais em experiências imersivas ou espacialmente contextualizadas. Esses ambientes podem apresentar persistência, interoperabilidade entre espaços, mecanismos de criação e circulação de ativos digitais e formas de atividade econômica, sendo sustentados por um conjunto variável de tecnologias, como realidade estendida, inteligência artificial, sistemas de comunicação, computação distribuída e tecnologias descentralizadas.}
\end{quote}

Essa definição preserva os elementos centrais da proposta inicial e incorpora a perspectiva tecnológica identificada na literatura. A palavra ``podem'' é deliberada: persistência, interoperabilidade e economia digital são características relevantes para a definição operacional deste trabalho, mas não devem ser tratadas como requisitos tecnológicos obrigatórios de toda implementação existente.

\section{Concordância com a unificação}

A unificação é considerada adequada para o objetivo deste projeto porque resolve uma diferença de escala entre as abordagens comparadas. A definição inicial descreve principalmente a experiência e as propriedades desejadas; Tukur et al. descrevem os meios tecnológicos e os contextos de aplicação; Ritterbusch e Teichmann ajudam a delimitar as características conceituais recorrentes.

A concordância, contudo, é parcial no sentido metodológico: não se deve afirmar que a definição unificada representa a definição científica definitiva do metaverso. Ela representa uma definição operacional deste trabalho, construída a partir da literatura analisada.

Essa distinção fortalece o texto acadêmico porque permite apresentar uma posição clara sem transformar uma síntese específica em consenso inexistente na literatura.
\chapter{Resultados}
\label{cha:resultados}

\section{Resultados da busca bibliográfica}

A busca bibliográfica realizada em 29 de setembro de 2026 confirmou a existência de um conjunto relevante e relativamente recente de revisões sobre o metaverso, especialmente em torno de quatro eixos: definição conceitual, arquitetura, tecnologias e aplicações.

A literatura de Ritterbusch e Teichmann concentra-se na definição do conceito e identificou 28 definições ou descrições após o processo de seleção e buscas de citações. Tukur et al. concentraram-se nas técnicas, tecnologias e aplicações e chegaram a 30 estudos primários em sua revisão de escopo. Mais recentemente, Boztas et al. analisaram 103 estudos, dos quais 33 continham definições explícitas e 25 apresentavam modelos arquiteturais. Esses resultados mostram que a discussão acadêmica já ultrapassou a fase de simplesmente perguntar o que é o metaverso e passou a discutir também como seus ambientes podem ser estruturados e implementados \cite{ritterbusch2023,tukur2024,boztas2025}.

Zhou, Chen e Jin acrescentam uma perspectiva sociotécnica ao relacionar definições, características técnicas e comportamento dos usuários. Bénaben, Congès e Fertier propõem uma definição apoiada na evolução das tecnologias imersivas e em um framework de propriedades. Enamorado-Díaz et al. analisam aplicações, tecnologias e desafios de engenharia de software \cite{zhou2024,benaben2025,enamorado2025}.

\section{Síntese temática}

A comparação dos trabalhos permite organizar a literatura em três níveis:

\begin{enumerate}
    \item \textbf{Nível conceitual:} o que caracteriza um metaverso, incluindo ambientes digitais compartilhados, interação, identidade, imersão e persistência;
    \item \textbf{Nível tecnológico:} tecnologias e infraestruturas que podem viabilizar essas características, incluindo XR, inteligência artificial, computação distribuída e tecnologias descentralizadas;
    \item \textbf{Nível aplicacional:} usos em educação, saúde, manufatura, negócios, entretenimento e outros setores.
\end{enumerate}

Essa organização sustenta a principal decisão conceitual do trabalho: tecnologias não devem ser automaticamente convertidas em requisitos da definição. A realidade virtual, por exemplo, pode intensificar a imersão, mas não precisa ser tratada como condição suficiente ou necessária para toda manifestação do metaverso.

\section{Comparação com a definição do projeto}

A definição originalmente adotada no projeto apresenta forte convergência com a literatura ao considerar ambientes tridimensionais compartilhados, identidades digitais, interação, imersão e persistência. A interoperabilidade também aparece como dimensão relevante em trabalhos recentes, embora permaneça associada a desafios técnicos e de padronização.

A principal contribuição da revisão é, portanto, reorganizar a definição em camadas. A camada conceitual descreve o fenômeno; a camada tecnológica explica os mecanismos que podem sustentá-lo; e a camada aplicacional descreve os contextos nos quais o ecossistema pode ser utilizado.

\section{Limitação dos resultados quantitativos}

Os números publicados por outras revisões foram utilizados apenas como referência metodológica. O presente trabalho não atribui a si esses quantitativos.

O fluxograma PRISMA numérico desta pesquisa permanece dependente da execução das strings diretamente nas interfaces das bases selecionadas e do registro dos resultados retornados em cada uma delas. Essa limitação é explicitada para preservar a validade do artigo e evitar que resultados de estudos anteriores sejam apresentados como se fossem dados coletados nesta revisão.

\chapter{Conclusão}
\label{cha:conclusao}

Este trabalho analisou a literatura científica sobre o metaverso a partir de diferentes perspectivas: definição, tecnologias, arquiteturas, aplicações e interoperabilidade. A análise foi deliberadamente limitada a informações que puderam ser verificadas nas próprias publicações utilizadas como fonte.

Os resultados mostram que não existe uma definição única consolidada. Ritterbusch e Teichmann analisaram 28 definições e descrições; Boztas et al. analisaram 103 estudos, dos quais 33 apresentavam definições explícitas e 25 apresentavam modelos arquiteturais; Tukur et al. analisaram técnicas, tecnologias e aplicações a partir de 30 estudos primários; e Enamorado-Díaz et al. selecionaram 35 estudos primários sobre aplicações e desafios de engenharia de software. \cite{ritterbusch2023,boztas2025,tukur2024,enamorado2025}

A comparação permite sustentar uma definição operacional em camadas. O primeiro nível descreve o fenômeno e suas propriedades; o segundo descreve as tecnologias que podem sustentá-lo; e o terceiro descreve suas aplicações. Essa estrutura reduz o risco de definir o metaverso por uma tecnologia específica.

A principal correção metodológica em relação às versões anteriores foi retirar números que não correspondiam a uma busca própria documentada. O presente trabalho não apresenta um PRISMA numérico próprio porque não dispõe de uma cadeia completa e auditável de identificação, duplicação, triagem e elegibilidade executada pelo estudo. Em seu lugar, os quantitativos das revisões analisadas são apresentados com sua fonte explícita.

Essa escolha é importante para a integridade acadêmica do trabalho. PRISMA 2020 orienta o relato transparente de revisões sistemáticas, enquanto PRISMA-S detalha a documentação das buscas para que possam ser reproduzidas. \cite{page2021,rethlefsen2021}

A definição operacional proposta é, portanto, uma síntese para este projeto, e não uma afirmação de consenso científico. Uma etapa futura poderá transformar o estudo em uma revisão sistemática própria, desde que sejam executadas as strings nas bases selecionadas, registrados os resultados de cada base, exportados os registros, removidas as duplicatas, documentados os motivos de exclusão e construído o fluxograma PRISMA a partir desses dados.

Até que essa etapa seja efetivamente realizada, qualquer número adicional de registros da revisão própria deve ser considerado inexistente, e não estimado.


\backmatter
\nocite{*}
\referencepage

\end{document}
